\section{Task Adaptation via Integrating Task-Specific Heuristics on Top of \plugmem}
\label{app:task_adaptation}
\subsection{LongMemEval}
\paragraph{Reflective Memory Management.}

RMM~\cite{tan2025prospectretrospectreflectivememory} introduces a novel reflective mechanism for long-term memory management. In general, RMM employs Prospective Reflection for knowledge extraction and memory updating, and Retrospective Reflection to refine memory retrieval in an online reinforcement learning manner.



Inspired by these ideas, we perform task adaptation of \plugmem to LongMemEval by incorporating the Prospective Reflection mechanism from RMM for semantic memory extraction and semantic node updating. Specifically, for memory extraction, we adopt the memory extraction prompt proposed in RMM to distill semantic memory for \plugmem. Compared to the original semantic memory extraction prompt used in \plugmem (Listing~\ref{lst:apdx_prompt_get_semantic}), the task-adapted version places greater emphasis on extracting user-specific personal information and includes in-context examples to better align with the LongMemEval setting. For memory updating, we follow the reflective updating strategy introduced in RMM to update semantic memory within the memory graph. Specifically, unlike the original \plugmem, which directly inserts newly extracted memory, the task-adapted version first retrieves the Top-K most similar existing semantic memories from the memory graph by computing cosine similarity between semantic memory embeddings and ranking them accordingly. The LLM is then queried to decide whether the new memory should be added as a new node or merged with an existing similar node. The prompt to extract memory and update memory is shown in Listing~\ref{lst:apdx_prompt_adpt_extract} and Listing~\ref{lst:apdx_prompt_adpt_update}.

\paragraph{Experiment.}

We conduct experiments on the multi-session subset of LongMemEval, with results shown in Table~\ref{tab:apdx_longmemeval_task_adpt}. The multi-session subset is among the most challenging portions of LongMemEval, as it requires the agent to answer questions that primarily involve counting the occurrences of a series of similar events across multiple sessions. This setting not only places significant demands on the retrieval module of the memory module to retrieve multiple gold memories simultaneously, but also challenges the structuring and reasoning module to generate precise semantic memories and perform accurate counting over them. The task-adapted version outperforms the original method for two main reasons: first, the task-adapted memory extraction prompt produces more precise semantic memories; second, the reflective updating strategy prevents redundant and repetitive semantic memories from being inserted into the memory graph, thereby enabling more efficient retrieval.

\begin{table}[ht]
\centering
\caption{Task Adaptation on the multi-session subset of LongMemEval.}
\label{tab:apdx_longmemeval_task_adpt}
\small
\setlength{\tabcolsep}{4pt}
\begin{tabular}{l|cccc}
\toprule
Method & Acc. (Multi-session) & PMI & $\text{\#Tok}_{\text{Avg.}}$ & Info. Density\\
\midrule
\plugmem & 64.7 & 5.06 & 311.9 & 1.62e-2 \\
\plugmem$_{\text{adapt}}$ & 70.7 & 5.62 & 327.5 & 1.71e-2\\
\bottomrule
\end{tabular}
\end{table}
\begin{lstlisting}[
  caption={Prompt template to extract semantic memory for task adaptation.},
  label={lst:apdx_prompt_adpt_extract},
]

============================================================
Prompt Extract_Memory
============================================================

Task Description: Given a session of dialogue between User and Agent, extract the personal summaries of User. Ensure the output adheres to the following rules:
Output results in OUTPUT format. The top-level tittle is "### Memory". The value should be a list of dictionaries, where each dictionary has the key "Summary":
- summary: A concise personal summary, which captures relevant information about User experiences, preferences, and background, across multiple turns.
If no personal summary can be extracted, return NO_TRAIT.
Example:
INPUT:
Turn 0:
- User: Did you check out that new gym in town?
- Agent: Yeah, I did. I am not sure I like the vibe there, though.
Turn 1:
- User: What was wrong with it?
- Agent: The folks there seemed to care more about how they looked than working out. It was a little too trendy for me. I am pretty plain.
Turn 2:
- User: Ah, got it. Well, maybe one of the older gyms will work out better for you-or I guess you could get that treadmill you were talking about before. Are you leaning one way or the other yet?
- Agent: I am leaning towards the treadmill. I think it will work better for my lifestyle. I just do not know which type to get. There are so many choices out there. Do you use a treadmill at your gym? Do you have a suggestion for a home one?
Turn 3:
- User: I usually just lift weights there, to be honest. But I think I have heard good things about the NordicTrack?
- Agent: Yeah, I have heard good things about that, too. I like the idea of a multi-exercise piece of equipment. As long as the weather is not too bad, then I prefer to go for a run. But since it rains quite a bit here, I like the idea of an inside option. How is the weather in New England?
OUTPUT:
### Memory:
1. **Summary:** User asked about a new gym in town and suggested older gyms or a treadmill as alternatives.
2. **Summary:** User usually lifts weights at the gym rather than using a treadmill.
3. **Summary:** User has heard good things about the NordicTrack treadmill.
Task: Follow the OUTPUT format demonstrated in the example above and extract the personal summaries for User from the following dialogue session.
Input: {episodic_memory}
Output:

\end{lstlisting}
\begin{lstlisting}[
  caption={Prompt template to update semantic memory for task adaptation.},
  label={lst:apdx_prompt_adpt_update},
]

============================================================
Prompt Update_Memory
============================================================

Task Description: Given a list of history personal summaries for a specific user and a new and similar personal summary from the same user, update the personal history summaries following the instructions below:
Input format: Both the history personal summaries and the new personal summary are provided in JSON format, with the top-level keys of "history_summaries" and "new_summary".
Possible update actions:
- Add: If the new personal summary is not relevant to any history personal summary, add it.
Format: Add()
- Merge: If the new personal summary is relevant to a history personal summary, merge them as an updated summary.
Format: Merge(index, merged_summary)
Note: index is the position of the relevant history summary in the list. merged_summary is the merged summary of the new summary and the relevant history summary. Two summaries are considered relevant if they discuss the same aspect of the user personal information or experiences.
Do not include additional explanations or examples in the output-only return the required action functions.
Example:
INPUT:
History Personal Summaries:
- {{"history_summaries": ["SPEAKER_1 works out although he does not particularlyenjoy it."]}}
New Personal Summary:
- {{"new_summary": "SPEAKER_1 exercises every Monday and Thursday."}}
OUTPUT ACTION:
Merge(0, SPEAKER_1 exercises every Monday and Thursday, although he does not particularly enjoy it.)
Task: Follow the example format above to update the personal history for the given case.
INPUT:
History Personal Summaries: {history_summaries}
New Personal Summary: {new_summary}
OUTPUT ACTION:

\end{lstlisting}

\subsection{HotpotQA}

\paragraph{Inspiration from HippoRAG2.}
Our retrieval design is partly inspired by HippoRAG2, which demonstrates that introducing an explicit graph structure and a controlled multi-hop traversal can outperform purely similarity-based retrieval on multi-hop QA.
In particular, HippoRAG2 separates \emph{recognition} (filtering/selecting relevant structured units) from \emph{expansion} (graph-based propagation to retrieve supporting evidence), reducing spurious hops and improving evidence quality.
\plugmem adopts this high-level principle but generalizes it to typed agent memory.
We construct semantic and procedural memory graphs and use abstract nodes (concepts/intents) as routing signals to activate specific proposition/prescription nodes, enabling abstraction-to-specificity traversal.
Moreover, analogous to HippoRAG2's recognition step, PlugMem employs an LLM-based controller to \textit{i)} select memory types (\textsc{GetMode}), \textit{ii)} propose next-hop abstract signals (\textsc{GetPlan}), and \textit{iii)} decide early stopping or focus candidates under a fixed budget (\textsc{MultiHopCtrl}).
Finally, we integrate a two-channel candidate update---link-based expansion through abstract nodes and embedding-based retrieval from a refined query---followed by reranking/pruning and optional reasoning-based compression. 

\paragraph{Test-time Scaling. }
Beyond the default configuration, \plugmem admits two practical test-time ``knobs'' that consistently improve performance.
First, we can increase the hop limit $T_{\max}$, effectively scaling the amount of retrieval computation at inference time.
This allows the retriever to chain additional abstraction--specificity steps and is analogous to test-time scaling in reasoning-centric methods. We experiment with a subset on HotpotQA, when increasing the multi-hop limit $T_{\max}$ from 2 to 4, we observe a performance gain from 66.00/74.73 to 69.00/78.11 (measured in Exact Match/F1-Score).


Secondly, when forming the next-hop query $Q_{t+1}$, we optionally replace naive concatenation of $(Q_t, \text{selected facts})$ with an LLM-based query synthesizer that rewrites a focused and self-contained query. Empirically, synthesis-based refinement reduces lexical drift and redundancy, and helps the embedding channel retrieve more targeted evidence in later hops.
We treat both mechanisms as controllable components within the retrieval controller, enabling a smooth accuracy--cost trade-off depending on the task budget.


\subsection{WebArena}
Prior work designed specifically for WebArena often incorporates manually engineered
\emph{domain-level instructions} (or ``tips'') into the agent’s input context.
A common pattern in task-specific adaptations—such as AWM~\cite{wang2024agentworkflowmemory},
SteP~\cite{sodhi2024stepstackedllmpolicies}, and more recently ColorBrowserAgent~\cite{zhou2026colorbrowseragentintelligentguiagent}—
is the inclusion of explicit guidance that biases the agent toward particular interaction
patterns within a site domain (e.g., navigation conventions or preferred UI elements).
In contrast, in our main experiments we deliberately avoid injecting such handcrafted
instructions or explicitly defined action policies.
Instead, we rely solely on \plugmem’s ability to automatically store, retrieve, and
re-apply both agent-generated and human demonstration trajectories as structured memory,
allowing behavior to emerge from experience rather than prompt engineering.

\paragraph{Further task adaptation}
We additionally conduct a targeted case study to examine whether \plugmem can benefit
from manual task adaptation when such domain-specific constraints are known.
We identify a subset of WebArena tasks that remain challenging for \plugmem and introduce
a small set of manually designed instructions into the reasoning module.
Specifically, we observe that \plugmem occasionally instructs the agent to navigate via
category dropdown menus when the correct strategy in WebArena is to use the search bar.
By incorporating the additional guidance shown in
Prompt~\ref{lst:webarena_task_adapt}, we resolve several previously failing cases
(see Example~\ref{tab:webarena_task_adapt_case}) by explicitly conditioning the agent’s
action selection. WebArena uses strict, programmatic success criteria but concise queries, where
otherwise reasonable strategies may be scored as incorrect.
Although \plugmem successfully retrieves relevant procedural knowledge for tasks such as
``searching for teeth grinding mouth guards,'' it does not, by default, infer that certain
plausible actions are disallowed by the benchmark’s implicit rules.
This case study demonstrates that \plugmem natively supports fine-grained manual
conditioning when necessary, while remaining effective without task-specific prompting
in the general case.
 
\begin{lstlisting}[
  caption=Additional instructions for the agent and reasonig module for further task-adaptation,
  label=lst:webarena_task_adapt,
]
============================================================
Prompt WebArena Task-specific Instructions
============================================================

1. When asked about a kind of product (not a category), search for the product using the search bar instead of using the category dropdown menu.
2. When asked about a category of product, access the dropdown menu through hover actions.
\end{lstlisting}
\begin{table}[t]
\centering
\small
\caption{Example of \plugmem with manual instructions.}
\label{tab:webarena_task_adapt_case}

\setlength{\tabcolsep}{8pt}
\renewcommand{\arraystretch}{1.1}
\begin{tabularx}{\linewidth}{p{0.97\linewidth}}

\toprule
\textbf{WebArena Manual Instruction Effect}\\
\midrule
\textbf{Query.} \emph{``\;What is the price range of teeth grinding mouth guard in the One Stop Market?\;''}\\
\midrule
\textbf{\plugmem with Manual Instruction.}\\
\(\triangleright\) \textbf{Reasoning}\\
The user's question is about determining the price range of wireless earphones in the One Stop Market. According to the guidelines, when asked about a specific kind of product (e.g., "wireless earphone"), the agent must use the search bar to locate the product instead of navigating through dropdown menus. The current observation shows that the search bar is present but appears disabled. However, based on the relevant experience provided, attempting to use the search bar is still recommended, as its functionality may not be accurately represented by its appearance.\\
\(\triangleright\) \textbf{Final Information}\\
1. Use the search bar to type "wireless earphone" and press Enter to locate relevant products.\\
2. Once the search results are displayed, review the prices of the listed wireless earphones to determine the price range.\\
3. If the search bar does not function, report the issue and consider alternative methods to locate the product, such as contacting support or exploring other navigation options.\\
\midrule
\textbf{AgentOccam Action 1:}
\textbf{type [1913] [wireless earphone] [1]}. (agent uses the search bar)\\
\midrule
...\\
\midrule
\textbf{\plugmem without Manual Instruction.}\\
\(\triangleright\) \textbf{Reasoning}\\
The user's question is about finding the price range of teeth grinding mouth guards in the One Stop Market. The observation does not provide any direct information about teeth grinding mouth guards or their prices. However, the relevant experience suggests that the agent should use the hierarchical menu structure to navigate to the appropriate category (likely under "Beauty \& Personal Care" or "Health \& Household") and then refine the search to locate the specific product type. If a price filter is available, it can be applied to determine the price range.\\
\(\triangleright\) \textbf{Final Information}\\
To answer the question, the agent should:\\
1. Hover over the hierarchical menu structure, starting with categories like "Beauty \& Personal Care" or "Health \& Household."\\
2. Look for subcategories related to oral care or dental care.\\
3. Navigate to the relevant subcategory (e.g., "Oral Care" or "Mouth Guards").\\
4. Review the products listed in the subcategory to identify teeth grinding mouth guards and note their prices.\\
5. If a price filter is available, apply it to determine the price range more efficiently.\\
6. Summarize the price range of the identified products.\\
\midrule
\textbf{AgentOccam Action 1:}\\
\textbf{hover [1534]}. (agent operates the dropdown menu (WRONG!))\\
\midrule
...\\
\bottomrule
\end{tabularx}
\end{table} 

